\documentclass[11pt]{article}

\usepackage[margin=1.08in]{geometry}
\usepackage{amsmath,amssymb,amsthm}
\usepackage{booktabs}
\usepackage[hidelinks]{hyperref}

\newtheorem{theorem}{Theorem}[section]
\newtheorem{proposition}[theorem]{Proposition}
\newtheorem{lemma}[theorem]{Lemma}
\newtheorem{corollary}[theorem]{Corollary}
\theoremstyle{definition}
\newtheorem{definition}[theorem]{Definition}
\theoremstyle{remark}
\newtheorem{remark}[theorem]{Remark}

\newcommand{\A}{\mathbb A}
\newcommand{\C}{\mathbb C}
\newcommand{\Q}{\mathbb Q}
\newcommand{\Z}{\mathbb Z}
\newcommand{\Gm}{\mathbb G_m}
\newcommand{\Spec}{\operatorname{Spec}}
\newcommand{\Frac}{\operatorname{Frac}}
\newcommand{\Cl}{\operatorname{Cl}}
\newcommand{\divisor}{\operatorname{div}}
\newcommand{\Supp}{\operatorname{Supp}}
\newcommand{\Jac}{\operatorname{Jac}}
\newcommand{\rank}{\operatorname{rank}}

\title{No Cubic \'Etale Morphism from the Pseudoplane
\(S(2,2,1)\) to the Affine Plane}
\author{Atharva Vaidya}
\date{25 July 2026}

\begin{document}
\maketitle

\begin{abstract}
Let
\[
 S=S(2,2,1)=
 \Spec\C[u,v,w]/\bigl(w^2-u-u^2v\bigr).
\]
We prove that there is no \'etale morphism
\(S\to\A^2_\C\) of geometric degree three.  The proof passes to the
finite normalization \(Y\to\A^2\) in the induced cubic function-field
extension.  Rational acyclicity of \(S\) forces every resolution boundary
to be a rational tree.  It follows that the singularities of \(Y\) have
rational-homology-sphere links and that the reduced boundary
\(Y\setminus S\) has vanishing compactly supported first cohomology.
An Euler and local-monodromy calculation then leaves one branch curve,
one boundary curve, and only \(2+1\) branch fibers.  The retained simple
sheet is consequently a homology line on \(S\).  Zaidenberg's
classification makes this line smooth, the
Abhyankar--Moh--Suzuki theorem rectifies its image in \(\A^2\), and the
resulting cyclic complement group cannot support connected cubic
monodromy generated by a transposition.

As a limited consequence, the result excludes geometric-degree-six
plane Keller maps that factor through the canonical quadratic chart
\(\A^2\to S(2,2,1)\) with cubic second stage.  It does not exclude
arbitrary degree-six Keller maps and does not prove the plane Jacobian
conjecture.
\end{abstract}

\section{Introduction and statement}

Put
\[
 B=\C[u,v,w]/(w^2-u-u^2v),\qquad S=\Spec B.
\]
The surface \(S\) is a smooth affine \(\Q\)-homology plane with
\(\pi_1(S)\simeq\Z/2\).  It is the \(\mathrm{ML}_1\) pseudoplane
\(S(2,2,1)\) in the notation of
Dubouloz--Palka~\cite{DuboulozPalka}; in particular it is not
\(\A^2\).

For a dominant morphism \(F:S\to\A^2\), its \emph{geometric degree} is
\[
 \deg_{\mathrm{geo}}F
 =[\C(S):\C(F)].
\]
Our main result is the following.

\begin{theorem}\label{thm:main}
There is no \'etale morphism
\[
 F:S(2,2,1)\longrightarrow\A^2_\C
\]
of geometric degree \(3\).
\end{theorem}

The genuinely cubic step is Theorem~\ref{thm:main}.  For context, the
same source admits no \'etale morphism to \(\A^2\) of degree one or two
either; see Corollary~\ref{cor:low-degree}.  The target matters.
Dubouloz and Palka construct many nonproper \'etale
\emph{endomorphisms} of \(S(2,2,1)\)~\cite{DuboulozPalka}.  Those maps
have target \(S(2,2,1)\), not \(\A^2\), and do not conflict with
Theorem~\ref{thm:main}.

The closest general nonexistence result located in the primary
literature is Miyanishi's theorem excluding \'etale morphisms
\(X\to\A^2\) when \(X\) is an \(\mathrm{ML}_0\) affine pseudoplane
(and also for Platonic \(\A^1\)-fiber spaces)
~\cite[Theorem 2.9]{Miyanishi}.  The present surface is
\(\mathrm{ML}_1\), so that theorem does not apply.  Miyanishi's later
lecture notes exclude maps from affine pseudoplanes of type
\((d,n,r)\) when \(r\ne2\)
~\cite[Theorem 2.4.10]{MiyanishiLectures}.  In the standard
\(\C^*\)-pseudoplane presentation, \(S(2,2,1)\) has \(r=2\), precisely
the case left open by that canonical-divisor argument.  To the
author's knowledge after a targeted review of these sources, the
degree-three exclusion below is not stated in the existing literature.
This is a cautious literature statement, not a claim of priority.

\section{The source and its finite normalization}

\subsection{Elementary properties of the pseudoplane}

\begin{lemma}\label{lem:source}
The surface \(S\) is smooth and rational, and
\[
 B^\times=\C^\times,\qquad
 H_i(S;\Q)=0\quad(i>0),\qquad \chi(S)=1.
\]
\end{lemma}

\begin{proof}
For \(h=w^2-u-u^2v\),
\[
 h_u=-1-2uv,\qquad h_v=-u^2,\qquad h_w=2w.
\]
These derivatives have no common zero, so \(S\) is smooth.  Moreover
\[
 \C(S)=\C(u,w),\qquad v=\frac{w^2-u}{u^2},
\]
so \(S\) is rational.

After inverting \(u\), the equation eliminates \(v\) and gives
\[
 B_u\simeq\C[u,u^{-1},w].
\]
Thus a unit of \(B\) becomes \(cu^n\) in \(B_u\).  If \(n>0\), the
same equality in \(B\) would make \(u\) a unit; if \(n<0\), apply this
argument to the inverse.  But \(u\) is not a unit because
\(B/(u,w)\simeq\C[v]\).  Hence \(n=0\) and \(B^\times=\C^\times\).

The morphism \(u:S\to\A^1\) is trivial over \(\Gm\):
\[
 S_{u\ne0}\simeq\Gm\times\A^1_w,
\]
while its reduced special fiber is \(V(u,w)\simeq\A^1_v\).  Additivity
of the compactly supported Euler characteristic yields
\[
 \chi(S)=\chi(\Gm)\chi(\A^1)+\chi(\A^1)=1.
\]
The standard topology of the pseudoplane gives
\(\pi_1(S)=\Z/2\); see, for example,
\cite[\S2 and Example \(S(k,r,a)\)]{DuboulozPalka}.
Therefore \(H_1(S;\Q)=0\).  A smooth affine complex surface has the
homotopy type of a finite CW complex of real dimension at most two
~\cite{Hamm}.  Hence its rational homology vanishes above degree two,
and the Euler equality forces \(H_2(S;\Q)=0\).
\end{proof}

\subsection{The normalization and generic branch type}

Assume for contradiction that \(F:S\to\A^2\) is \'etale of geometric
degree three.  Write
\[
 K=\Frac B,\qquad L=\C(F),\qquad A=\C[F]\simeq\C[x,y].
\]
Let \(A'\) be the integral closure of \(A\) in \(K\), and set
\[
 \pi:Y=\Spec A'\longrightarrow\Spec A=\A^2.
\]

\begin{lemma}\label{lem:normalization}
The map \(\pi\) is finite flat of rank three.  There is a canonical open
immersion \(S\hookrightarrow Y\) through which \(F\) factors, and
\(\pi|_S=F\).
\end{lemma}

\begin{proof}
The assumption on geometric degree makes \(F\) dominant, so its two
coordinates are algebraically independent and
\(A\simeq\C[x,y]\).  Finiteness follows from excellence of \(A\).
Every element of \(A'\)
is integral over \(A\), hence integral over \(B\); since \(B\) is
normal, \(A'\subset B\).  The induced map \(S\to Y\) is birational and
quasi-finite: each of its fibers lies in a finite fiber of the
quasi-finite morphism \(F\).  The birational form of Zariski's Main
Theorem makes it an open immersion.

Every two-dimensional normal local ring is Cohen--Macaulay.  After
localizing over a point of \(\Spec A\), a parameter system of \(A\)
remains a parameter system in every local ring of the finite
\(A\)-algebra \(A'\).  Cohen--Macaulayness therefore makes \(A'\) a
maximal Cohen--Macaulay module over the regular local ring \(A\).
Auslander--Buchsbaum, equivalently miracle flatness, makes it free.
Hence \(A'\) is locally free, and its generic rank is three.
\end{proof}

Let
\[
 R=(Y\setminus S)_{\mathrm{red}}
\]
and let \(\Delta\subset\A^2\) be the reduced branch curve of \(\pi\).
The branch curve is nonempty.  Indeed, \(Y\) is normal and integral,
\(\A^2\) is regular, and \(\pi\) is finite and generically separable.
Thus purity of the branch locus would otherwise make the connected
finite cover \(Y\to\A^2\) \'etale, contrary to the triviality of
\(\pi_1^{\mathrm{et}}(\A^2_\C)\).

\begin{lemma}\label{lem:branch-type}
At the generic point of every irreducible component
\(\Delta_j=V(f_j)\) of \(\Delta\),
\[
 \divisor_Y(f_j)=2E_j+\overline C_j,
 \tag{2.1}
\]
where \(E_j\subset R\) is ramified and
\(\overline C_j\) is unramified, both with residue degree one.  In
particular, if \(r\) and \(c\) denote the numbers of irreducible
components of \(R\) and \(\Delta\), respectively, then
\[
 r\ge c\ge1.
 \tag{2.2}
\]
\end{lemma}

\begin{proof}
Total cubic ramification above \(\Delta_j\) would leave one prime
\(E_j\), of ramification index three.  Since \(\pi|_S\) is \'etale,
that entire prime would lie outside \(S\).  The principal inverse image
of \(\Delta_j\) has no additional zero-dimensional component:
\(Y\) is Cohen--Macaulay and \(f_j\) is a nonzerodivisor, so
\(\mathcal O_Y/(f_j)\) is Cohen--Macaulay of pure dimension one, with
no embedded or isolated component.  Finiteness makes every curve
component dominate \(\Delta_j\).  Thus
\(F^{-1}(\Delta_j)=\varnothing\), making
\[
 f_j(F)\in B^\times=\C^\times.
\]
This is impossible because \(A=\C[F]\hookrightarrow B\) and \(f_j\)
is nonconstant.

Over \(\C\) the cubic extension is tame.  Once total ramification is
excluded, the degree formula over the generic discrete valuation ring
has only the decomposition
\[
 3=2\cdot1+1\cdot1,
\]
which gives (2.1).  The ramified prime is omitted from \(S\), while the
unramified prime must meet \(S\), for otherwise the same unit argument
would apply.  Distinct branch components have distinct ramified primes,
because the finite image of an irreducible curve is irreducible.  This
proves (2.2).
\end{proof}

\begin{lemma}[Purity of the boundary]\label{lem:boundary-pure}
The closed set \(Y\setminus S\) has no isolated points.  Consequently
its underlying set is the pure curve \(R\).
\end{lemma}

\begin{proof}
Suppose that \(z\in Y\setminus S\) were isolated away from the
divisorial boundary.  Since \(Y\) is affine and has only finitely many
boundary curves, choose from each boundary prime an element not
vanishing at \(z\), and take their product.  This gives
\(g\in\mathcal O(Y)\) vanishing on every boundary curve but not at
\(z\).  Then
\[
 W=D_Y(g)
\]
is a normal affine surface and
\[
 Z=W\setminus(S\cap W)
\]
is finite and nonempty.  On the other hand,
\(S\cap W=D_S(g)\) is affine.  Normality and Hartogs extension across
codimension two give
\[
 \Gamma(W\setminus Z,\mathcal O)=\Gamma(W,\mathcal O).
\]
If \(W\setminus Z\) were affine, its canonical affinization would
identify the open immersion \(W\setminus Z\hookrightarrow W\) with an
isomorphism, contradicting \(Z\ne\varnothing\).  Thus no such \(z\)
exists.
\end{proof}

\section{The log-topological collapse}

This section proves the global input that makes the cubic argument
rigid.  No classification of cubic covers is used.

\begin{proposition}[Rational-tree boundary]\label{prop:rational-tree}
For every smooth projective simple-normal-crossing completion
\((X,B_X)\) of \(S\), every component of \(B_X\) is rational and its
dual graph is a tree.
\end{proposition}

\begin{proof}
The surface \(X\) is a smooth projective model of the rational function
field \(\C(u,w)\), hence is rational and
\(H^1(X;\Q)=0\).  Relative cohomology identifies
\[
 H^i(X,B_X;\Q)\simeq H_c^i(S;\Q).
\]
Poincar\'e duality on \(S\), together with
Lemma~\ref{lem:source}, gives
\[
 H_c^1(S;\Q)\simeq H_3(S;\Q)=0,\qquad
 H_c^2(S;\Q)\simeq H_2(S;\Q)=0.
\]
The long exact sequence first makes \(B_X\) connected and then gives
\(H^1(B_X;\Q)=0\).  For an SNC projective curve,
\[
 \dim H^1(B_X;\Q)
 =2\sum_i g(B_{X,i})+b_1(\Gamma_{B_X}),
\]
where \(\Gamma_{B_X}\) is the dual graph.  Every term is nonnegative,
so all genera and the graph's first Betti number vanish.
\end{proof}

\begin{lemma}[Links and Alexander--Lefschetz duality]
\label{lem:alexander}
Let \(Z\) be a normal complex algebraic surface.  If the link of every
singular point is a rational homology three-sphere, then \(Z\), with
its complex orientation, is an oriented rational homology
four-manifold.  For every closed locally compact subset \(T\subset Z\),
cap-product duality gives
\[
 H_i(Z,Z\setminus T;\Q)\simeq H_c^{4-i}(T;\Q).
 \tag{3.0}
\]
\end{lemma}

\begin{proof}
The singular locus of a normal surface has complex codimension at
least two and is therefore discrete.  At a smooth point the local
homology is that of \(\R^4\).  At a singular point, the local conical
structure identifies the local homology with the shifted reduced
homology of the link.  The rational-homology-sphere hypothesis gives
the same local groups as \(\R^4\), and the complex orientation on the
smooth locus supplies the orientation class.  The displayed
isomorphism is Alexander--Lefschetz duality for an oriented homology
manifold, equivalently Verdier duality for its constant rational
sheaf; see \cite[Chapter~3]{Dimca}.
\end{proof}

\begin{proposition}\label{prop:boundary-cohomology}
Every singular point of \(Y\) has a rational-homology-sphere link, and
\[
 H_c^1(R;\Q)=0.
 \tag{3.1}
\]
Consequently every irreducible component of \(R\) has normalization
\(\A^1\), distinct components are disjoint, and every point of \(R\)
is analytically unibranch.  In particular,
\[
 \chi(R)=r.
 \tag{3.2}
\]
\end{proposition}

\begin{proof}
Resolve \(Y\) by a morphism that is an isomorphism over \(S\), complete
the resolution, and make the boundary SNC.  The result is a completion
as in Proposition~\ref{prop:rational-tree}.  The reduced exceptional
divisor over a singular point \(p\in Y\setminus S\) is connected by
the connectedness theorem for a resolution of a normal surface.  It is
a subdivisor of that rational tree and therefore itself a tree of
smooth rational curves; its intersection matrix is negative definite.
The plumbing calculation for a normal surface singularity
~\cite{Mumford} gives
\[
 b_1(\operatorname{Link}(p,Y);\Q)
 =2\sum_C g(C)+b_1(\Gamma)=0.
\]
Thus every link is a rational homology sphere, and
Lemma~\ref{lem:alexander} applies to \(Y\).

Both \(Y\) and \(S\) are affine surfaces.  Hamm's theorem gives
\(H_3(Y;\Q)=0\), while Lemma~\ref{lem:source} gives
\(H_2(S;\Q)=0\).  The pair sequence therefore yields
\[
 H_3(Y,S;\Q)=0.
\]
Equation (3.0), together with \(S=Y\setminus R\), identifies
\[
 H_3(Y,Y\setminus R;\Q)\simeq H_c^1(R;\Q),
\]
and proves (3.1).

For completeness, let
\[
 \nu:\widetilde R=\coprod_{i=1}^r\widetilde R_i\longrightarrow R
\]
be the normalization.  Let \(g_i\) be the genus of the smooth
projective completion of \(\widetilde R_i\), let \(s_i\ge1\) be its
number of punctures, and let \(b_p\) be the number of analytic branches
of \(R\) at \(p\).  The normalization exact sequence is
\[
0\longrightarrow\Q_R\longrightarrow\nu_*\Q_{\widetilde R}
\longrightarrow
\bigoplus_{p\in R}\Q_p^{\,b_p-1}\longrightarrow0.
\]
Every normalized affine component has at least one puncture.
Compactly supported cohomology gives
\[
 \dim H_c^1(R;\Q)
 =\sum_{i=1}^r(2g_i+s_i-1)
  +\sum_{p\in R}(b_p-1).
 \tag{3.3}
\]
All terms are nonnegative.  Their vanishing gives
\(g_i=0\), \(s_i=1\), and \(b_p=1\).  Hence
\(\widetilde R_i\simeq\A^1\), and the remaining assertions, including
(3.2), follow.
\end{proof}

\begin{lemma}[Local support and monodromy]\label{lem:local-cover}
Let \(\rho:Z\to W\) be a finite flat generically \'etale morphism
from a normal complex surface to a smooth complex surface, and let
\(D\subset W\) be its branch curve.
\begin{enumerate}
\item If \(z\in\rho^{-1}(q)\) has local fiber length one, then
  \(\rho\) is \'etale at \(z\).
\item The support points of \(\rho^{-1}(q)\) are in bijection with the
  orbits of the local monodromy image on the geometric generic fiber.
\item Every analytic branch of \(D\) at \(q\) has a ramification curve
  germ above it whose closure ends at a non-\'etale support point over
  \(q\).
\end{enumerate}
\end{lemma}

\begin{proof}
For (1), the local fiber algebra at \(z\) has length one and is
\(\C\).  Base change for relative differentials gives
\(\Omega_{Z/W,z}\otimes\C=0\); Nakayama's lemma gives
\(\Omega_{Z/W,z}=0\).  Flatness then makes \(\rho\) \'etale at \(z\).

For (2), take a sufficiently small analytic ball about \(q\).
Its inverse image separates into neighborhoods of the finitely many
points over \(q\).  Normality makes each surface germ irreducible, and
the complement of \(\rho^{-1}(D)\) in an irreducible normal surface
germ is connected.  Each such punctured neighborhood is therefore one
connected component of the unbranched cover over the ball minus
\(D\), hence one local-monodromy orbit; conversely every orbit has one
such closure.  This is also the local form of the comparison between
finite \'etale and finite topological covers
\cite[Expos\'e XII]{SGA1}.

For (3), restrict to a smooth point of the chosen branch near \(q\).
Its generic ramification prime exists by the definition of \(D\).
The closure of that prime meets the fiber over \(q\), because \(\rho\)
is finite, and the endpoint remains in the closed non-\'etale locus.
\end{proof}

\begin{proposition}\label{prop:branch-topology}
Every irreducible component of \(\Delta\) has normalization \(\A^1\).
Different components are disjoint, and every point of \(\Delta\) is
analytically unibranch.  Hence
\[
 \chi(\Delta)=c.
 \tag{3.4}
\]
\end{proposition}

\begin{proof}
Suppose that \(q\in\Delta\) had two distinct analytic branches.  Each
branch is generically simply ramified by Lemma~\ref{lem:branch-type}.
Lemma~\ref{lem:local-cover}(3) specializes its ramification lift to a
non-\'etale support point of the length-three fiber \(\pi^{-1}(q)\).
By Lemma~\ref{lem:local-cover}(1), every such support has length at
least two.  Two distinct support points would therefore require fiber
length at least four.

All ramification branches over the target branches must consequently
specialize to one point \(p\in R\).  If \(R\) had only one analytic
branch there, its finite analytic image would be irreducible and could
not contain two distinct target branches.  Thus distinct target
branches force distinct analytic branches of \(R\) at \(p\), contrary to
Proposition~\ref{prop:boundary-cohomology}.  Thus \(\Delta\) is
unibranch everywhere; in particular its irreducible components cannot
meet.

For a component \(\Delta_j\), the finite map
\(E_j\to\Delta_j\) in Lemma~\ref{lem:branch-type} is birational.
It identifies the normalizations.  Proposition
\ref{prop:boundary-cohomology} therefore gives
\[
 \widetilde\Delta_j\simeq\widetilde E_j\simeq\A^1.
\]
Because every point is unibranch, the normalization map is bijective
and hence a homeomorphism in the complex topology.  Thus each component
has Euler characteristic one, proving (3.4).
\end{proof}

\begin{proposition}[Cubic collapse]\label{prop:collapse}
The branch curve and the boundary are irreducible, and every affine
branch fiber has two support points:
\[
 r=c=1.
 \tag{3.5}
\]
Equivalently, no local inertia group at an affine branch point acts
transitively on all three sheets.
\end{proposition}

\begin{proof}
Choose a finite set \(\Sigma\subset\Delta\) containing every point at
which the support fiber differs from the generic simple-branch fiber
and every zero-dimensional constructible stratum of the finite map.
Let \(o_p\) denote the number of orbits of the local monodromy image
at \(p\) on the three sheets.  Lemma~\ref{lem:local-cover}(2) says
that the support fiber has \(o_p\) points.  Since \(p\in\Delta\), the
local image is nontrivial.  A generic simple branch fiber has
\(o_p=2\); at an exceptional branch point the possibilities are
\[
\begin{array}{c|c|c}
\text{local inertia}&\text{orbits}&o_p\\ \hline
\langle(12)\rangle&\{1,2\},\{3\}&2\\
\langle(123)\rangle&\{1,2,3\}&1\\
S_3&\{1,2,3\}&1.
\end{array}
\]
Let \(N_{\mathrm{tr}}\) count the points with \(o_p=1\).
After stratifying by the finite support-fiber cardinality,
constructible additivity of Euler characteristic gives
\[
\begin{aligned}
\chi(Y)
 &=3\chi(\A^2\setminus\Delta)
   +2\chi(\Delta\setminus\Sigma)+\sum_{p\in\Sigma}o_p\\
 &=3-\chi(\Delta)-N_{\mathrm{tr}}.
\end{aligned}
\tag{3.6}
\]
On the other hand,
\[
\chi(Y)=\chi(S)+\chi(R)=1+r.
\]
Using Proposition~\ref{prop:branch-topology} gives the exact balance
\[
 r+c+N_{\mathrm{tr}}=2.
\tag{3.7}
\]
Together with \(r\ge c\ge1\) from
Lemma~\ref{lem:branch-type}, this has the unique solution
\[
 r=c=1,\qquad N_{\mathrm{tr}}=0.
\]
\end{proof}

\section{The retained sheet is a smooth affine line}

Write \(\Delta\) for the unique branch curve, \(E\) for the unique
boundary curve, and \(\overline C\) for the retained prime.  Thus
generically
\[
 \divisor_Y(f)=2E+\overline C
\tag{4.1}
\]
for an irreducible equation \(f\) of \(\Delta\).

\begin{proposition}\label{prop:retained}
The two curves are disjoint, \(\overline C\cap E=\varnothing\).
Moreover \(\overline C\) is contained in \(S\), and restriction induces
an isomorphism
\[
 F|_C:C:=\overline C\xrightarrow{\ \sim\ }\Delta.
\tag{4.2}
\]
\end{proposition}

\begin{proof}
Both finite birational maps
\[
 E\longrightarrow\Delta,\qquad
 \overline C\longrightarrow\Delta
\]
induce isomorphisms between their normalizations.  Thus both
normalized curves identify with the common normalization
\(\A^1\to\Delta\).
Proposition~\ref{prop:branch-topology} says that this normalization is
bijective.  Hence each of \(E\) and \(\overline C\) has exactly one
point above every \(q\in\Delta\).

If \(p\in E\cap\overline C\) and \(q=\pi(p)\), these two unique
specializations would coincide.  There is no further point of
\(\pi^{-1}(q)\).  Indeed, the reduced support of the
Cohen--Macaulay divisor \(V(f)\) is exactly
\(E\cup\overline C\): purity excludes isolated components, finiteness
prevents a curve from mapping to \(q\), and (4.1) lists all divisorial
components.  The
length-three fiber would therefore have one support point, so its local
inertia would be transitive.  This contradicts
Proposition~\ref{prop:collapse}.  Thus
\[
 E\cap\overline C=\varnothing.
\]
Since Lemma~\ref{lem:boundary-pure} and (3.5) give
\(Y\setminus S=E\) set-theoretically, \(\overline C\subset S\).

The scheme-theoretic inverse image \(F^{-1}(\Delta)\) is reduced,
because it is the \'etale base change of the reduced curve \(\Delta\).
Its underlying set is \(C=\overline C\).  Thus
\[
 F|_C:C\longrightarrow\Delta
\]
is finite, surjective, \'etale, and generically of degree one.  A
finite \'etale algebra of rank one is the base algebra itself.  Here
finiteness comes from restricting \(\pi\), and the rank is constantly
one because \(\Delta\) is integral.  This proves (4.2).
\end{proof}

\begin{proposition}\label{prop:smooth-line}
The curves \(C\) and \(\Delta\) are smooth and isomorphic to \(\A^1\).
\end{proposition}

\begin{proof}
The finite bijective normalization
\(\A^1\to\Delta\) is a complex-topological homeomorphism.  By
Proposition~\ref{prop:retained}, \(C\) is therefore an irreducible
affine curve homeomorphic to \(\A^1(\C)\), that is, a homology line on
the smooth \(\Q\)-homology plane \(S\).

Zaidenberg's homology-line theorem states that if a homology line
\(\Gamma\) on a smooth \(\Q\)-homology plane \(X\) is singular, then
\(X\simeq\A^2\); more precisely the pair rectifies to a coprime
curve \(x^k-y^\ell=0\) in \(\A^2\)
~\cite[Theorem 1(c)]{Zaidenberg}.  Since
\(\pi_1(S)=\Z/2\), the surface \(S\) is not \(\A^2\).  Hence \(C\) is
smooth and \(C\simeq\A^1\).  The isomorphism (4.2) gives the same
conclusion for \(\Delta\).
\end{proof}

\begin{remark}
Smoothness is the exact conclusion needed here; the retained curve need
not be a fiber of the standard \(\A^1\)-fibration
\(u:S\to\A^1\).  For example
\[
 V(v)\subset S,\qquad
 B/(v)\simeq\C[u,w]/(w^2-u)\simeq\C[w],
\]
is a principal affine line on which \(u=w^2\) is nonconstant.
\end{remark}

\section{The monodromy contradiction}

\begin{proof}[Proof of Theorem~\ref{thm:main}]
By Proposition~\ref{prop:smooth-line}, \(\Delta\) is a smooth closed
embedding of \(\A^1\) in \(\A^2\).  The
Abhyankar--Moh--Suzuki theorem rectifies it
~\cite{AbhyankarMoh,Suzuki}.  Therefore
\[
 V:=\A^2\setminus\Delta\simeq\A^1\times\Gm,
\qquad \pi_1(V)\simeq\Z.
\tag{5.1}
\]
Equations (4.1) and (4.2) give
\[
 U:=S\setminus C
 =Y\setminus(E\cup C)
 =\pi^{-1}(V).
\]
Thus \(U\to V\) is a finite \'etale cover of degree three.  It is
connected because \(U\) is a nonempty open subset of the irreducible
surface \(S\).  Its monodromy image is consequently a transitive
subgroup of \(S_3\), by the comparison between finite \'etale and
finite topological covers over \(\C\)
\cite[Expos\'e XII]{SGA1}.

By (5.1), that image is cyclic and generated by a meridian around
\(\Delta\).  The only transitive cyclic subgroup of \(S_3\) is the
order-three subgroup generated by a three-cycle.  But (4.1) says that
the generic branch meridian acts as a transposition.  The subgroup
generated by a transposition has two orbits, a contradiction.
\end{proof}

\section{Low-degree and plane-map consequences}

\begin{proposition}\label{prop:galois}
Let \(X=\Spec R\) be a normal integral affine complex surface with
\(R^\times=\C^\times\).  If \(G:X\to\A^2\) is \'etale and
\(\C(X)/\C(G)\) is finite Galois, then the extension is trivial.
\end{proposition}

\begin{proof}
Form the finite normalization \(Z\to\A^2\) in the extension and embed
\(X\) as an open subset as in Lemma~\ref{lem:normalization}.  A
nontrivial connected normalization must have a branch divisor, by
purity and the trivial \'etale fundamental group of \(\A^2\).  Galois
transitivity makes every prime over a branch component ramified, hence
disjoint from the \'etale open set \(X\).  Pulling back an equation of
that component then gives a nonconstant unit of \(R\), a contradiction.
\end{proof}

\begin{corollary}\label{cor:low-degree}
There is no \'etale morphism \(S(2,2,1)\to\A^2\) of geometric degree
at most three.
\end{corollary}

\begin{proof}
Degree three is Theorem~\ref{thm:main}.  A separable degree-two
extension is Galois, so degree two follows from
Proposition~\ref{prop:galois}.

In degree one, Zariski's Main Theorem identifies \(S\) with an affine
open subset \(U\subset\A^2\).  If \(\A^2\setminus U\) contained a
curve, its equation would restrict to a nonconstant unit on \(S\).
Thus the complement is finite.  If it were nonempty, normal Hartogs
extension would give
\(\Gamma(U,\mathcal O)=\C[x,y]\); since \(U\) is affine, its
affinization would then fill the missing points.  Hence \(U=\A^2\),
contrary to \(\pi_1(S)=\Z/2\).
\end{proof}

There is a canonical quadratic \'etale chart
\[
 \pi_2:\A^2_{a,b}\longrightarrow S
\]
given by
\[
 u=a^2,\qquad
 w=a(1+2a^2b),\qquad
 v=4b(1+a^2b).
\tag{6.1}
\]
On \(B\), use the hypersurface Poisson bracket
\[
 \{u,v\}=-2w,\qquad
 \{u,w\}=-u^2,\qquad
 \{v,w\}=1+2uv.
\tag{6.2}
\]
Its coefficient vector
\((-2w,-u^2,1+2uv)\) never vanishes on \(S\), by the smoothness
calculation in Lemma~\ref{lem:source}.  Hence the induced bivector on
the smooth surface \(S\) is everywhere nondegenerate.  In particular,
\(\{P,Q\}\in\C^\times\) makes \((P,Q):S\to\A^2\) \'etale.
A direct calculation gives
\[
 \Jac_{a,b}(\pi_2^*P,\pi_2^*Q)
 =-4\,\pi_2^*\{P,Q\}.
\tag{6.3}
\]
Because the bracket is nondegenerate, this identity also proves that
\(\pi_2\) is \'etale.

\begin{corollary}[Limited plane Keller consequence]
\label{cor:degree-six}
There is no Darboux pair \(P,Q\in B\) satisfying
\[
 \{P,Q\}=1,\qquad [\Frac B:\C(P,Q)]=3.
\]
Equivalently, no plane Keller map
\[
 H=(\pi_2^*P,\pi_2^*Q):\A^2_{a,b}\longrightarrow\A^2
\]
arising through (6.1) has geometric degree six with a cubic second
stage.
\end{corollary}

\begin{proof}
The preceding nondegeneracy makes \((P,Q):S\to\A^2\) \'etale, so the first
assertion is Theorem~\ref{thm:main}.  Furthermore
\[
 \C(a,b)=\C(S)(a),\qquad a^2=u,
\]
because generically
\[
 b=\frac{w/a-1}{2u}.
\]
The element \(u\) is not a square in the rational function field
\(\C(u,w)\), as its \(u\)-adic valuation is one.  Thus this extension
is genuinely quadratic.  The tower law would therefore
give
\[
 [\C(a,b):\C(P,Q)]=2\cdot3=6,
\]
while (6.3) gives \(\Jac(H)=-4\).
\end{proof}

\begin{remark}[Logical boundary]\label{rem:scope}
Corollary~\ref{cor:degree-six} concerns only plane Keller maps that
factor through the fixed chart (6.1), equivalently through its
specified invariant subring \(B\).  It does not exclude arbitrary
plane Keller maps of geometric degree six.  Geometric degree is also
different from the maximum total degree of the two coordinate
polynomials.  Consequently this result supplies no new total-degree
bound for a hypothetical plane Jacobian counterexample and does not
resolve the two-dimensional Jacobian conjecture.
\end{remark}

\section{Verification, disclosure, and remaining review}

The proof is geometric.  Companion scripts check the elementary
identities and finite combinatorics but are not substitutes for the
theorem-level arguments.  From the repository root, run
\begin{verbatim}
.venv/bin/python \
  current_context/verify_route_a_log_topology_cubic_collapse.py
.venv/bin/python \
  current_context/verify_route_a_contractible_retained_sheet_cubic_exclusion.py
\end{verbatim}
The first script checks the rational-acyclic Betti calculation, the
compactly supported curve formula, the cubic fiber-length obstruction,
and the unique Euler-balance solution.  The second checks smoothness of
\(S\), the local length-three algebras, the rank-one \'etale step, and
the cyclic subgroup obstruction in \(S_3\).

OpenAI Codex was used in the research workflow for symbolic
verification, literature discovery, drafting, and adversarial proof
review.  The named author is responsible for every statement.  Before
submission, the proof should receive an independent human line-by-line
check, with particular attention to:
\begin{enumerate}
\item rational Alexander--Lefschetz duality for the singular
  normalization \(Y\);
\item the passage from local branch length to unibranchness in
  Proposition~\ref{prop:branch-topology};
\item the exact hypotheses of Zaidenberg's Theorem~1(c); and
\item a broader MathSciNet, zbMATH, and expert literature review before
  any priority statement.
\end{enumerate}

\begin{thebibliography}{99}

\bibitem{AbhyankarMoh}
S.~S. Abhyankar and T.~T. Moh,
\emph{Embeddings of the line in the plane},
J. Reine Angew. Math. \textbf{276} (1975), 148--166.
\href{https://doi.org/10.1515/crll.1975.276.148}
{doi:10.1515/crll.1975.276.148}.

\bibitem{Dimca}
A. Dimca,
\emph{Sheaves in Topology},
Universitext, Springer, Berlin, 2004.
\href{https://doi.org/10.1007/978-3-642-18868-8}
{doi:10.1007/978-3-642-18868-8}.

\bibitem{DuboulozPalka}
A. Dubouloz and K. Palka,
\emph{The Jacobian Conjecture fails for pseudo-planes},
Adv. Math. \textbf{339} (2018), 248--284.
\href{https://doi.org/10.1016/j.aim.2018.09.020}
{doi:10.1016/j.aim.2018.09.020};
\href{https://arxiv.org/abs/1701.01425}{arXiv:1701.01425}.

\bibitem{Hamm}
H.~A. Hamm,
\emph{Zum Homotopietyp Steinscher R\"aume},
J. Reine Angew. Math. \textbf{338} (1983), 121--135.
\href{https://doi.org/10.1515/crll.1983.338.121}
{doi:10.1515/crll.1983.338.121}.

\bibitem{Miyanishi}
M. Miyanishi,
\emph{Affine pseudo-coverings of algebraic surfaces},
J. Algebra \textbf{293} (2005), 156--176.
\href{https://doi.org/10.1016/j.jalgebra.2005.01.042}
{doi:10.1016/j.jalgebra.2005.01.042}.

\bibitem{Mumford}
D. Mumford,
\emph{The topology of normal singularities of an algebraic surface
and a criterion for simplicity},
Publ. Math. Inst. Hautes \'Etudes Sci. \textbf{9} (1961), 5--22.
\href{https://doi.org/10.1007/BF02698717}
{doi:10.1007/BF02698717}.

\bibitem{MiyanishiLectures}
M. Miyanishi,
\emph{Lectures on geometry and topology of polynomials:
surrounding the Jacobian conjecture},
lecture notes, 2015.
\href{https://arxiv.org/abs/1504.07179}{arXiv:1504.07179}.

\bibitem{SGA1}
A. Grothendieck,
\emph{Rev\^etements \'etales et groupe fondamental (SGA 1)},
Lecture Notes in Mathematics \textbf{224}, Springer, Berlin, 1971.
\href{https://doi.org/10.1007/BFb0058656}
{doi:10.1007/BFb0058656}.

\bibitem{Suzuki}
M. Suzuki,
\emph{Propri\'et\'es topologiques des polyn\^omes de deux variables
complexes, et automorphismes alg\'ebriques de l'espace \(\C^2\)},
J. Math. Soc. Japan \textbf{26} (1974), 241--257.
\href{https://doi.org/10.2969/jmsj/02620241}
{doi:10.2969/jmsj/02620241}.

\bibitem{Zaidenberg}
M. Zaidenberg,
\emph{Affine lines on \(\Q\)-homology planes and group actions},
Transform. Groups \textbf{11} (2006), 725--735.
\href{https://doi.org/10.1007/s00031-005-1122-5}
{doi:10.1007/s00031-005-1122-5};
\href{https://arxiv.org/abs/math/0611335}{arXiv:math/0611335}.

\end{thebibliography}

\end{document}
